\section{Evaluasi}
\label{sec:evaluasi-bab15}

\subsection{Soal Pilihan Ganda}

\begin{enumerate}
    \item Perintah DEBUG yang menampilkan seluruh register dan
          \eng{flag} adalah \dots
          \begin{enumerate}
              \item \code{U} \item \code{D}
              \item \code{R} \item \code{T}
          \end{enumerate}

    \item Pada saat \lstasm|CALL HITUNG| ditelusuri, perintah yang
          \emph{tidak} masuk ke dalam prosedur \asm{HITUNG} adalah
          \dots
          \begin{enumerate}
              \item \code{T} \item \code{P}
              \item \code{U} \item \code{R}
          \end{enumerate}

    \item Keluaran baris \eng{flag} DEBUG yang menunjukkan bahwa
          \flag{ZF} bernilai satu dan \flag{CF} bernilai nol adalah
          \dots
          \begin{enumerate}
              \item \code{NV UP EI PL NZ NA PO NC}
              \item \code{NV UP EI PL ZR NA PE NC}
              \item \code{OV DN DI NG NZ AC PO CY}
              \item \code{NV UP EI NG NZ NA PO CY}
          \end{enumerate}

    \item Sebuah prosedur menekan \reg{AX}, \reg{BX}, dan \reg{CX},
          tetapi hanya memulihkan \reg{CX} dan \reg{BX}. Akibat yang
          paling mungkin terjadi adalah \dots
          \begin{enumerate}
              \item assembler menolak program karena jumlah
                    \asm{PUSH} dan \asm{POP} tidak seimbang
              \item \asm{RET} melompat ke alamat yang bukan alamat
                    kembali
              \item register \reg{AX} otomatis dipulihkan oleh DOS
              \item program berjalan benar karena \reg{AX} tidak
                    dipakai lagi
          \end{enumerate}

    \item Program selesai tanpa pesan kesalahan dari assembler dan
          sistem operasi, tetapi hasil perhitungannya salah. Jenis
          kesalahan ini adalah \dots
          \begin{enumerate}
              \item kesalahan sintaks \item kesalahan logika
              \item kesalahan segmentasi \item \eng{overflow}
                    \eng{stack}
          \end{enumerate}
\end{enumerate}

\subsection{Soal Uraian}

\begin{enumerate}
    \item Jelaskan perbedaan perintah \code{T} dan \code{P} ketika
          dihadapkan pada instruksi \asm{CALL} dan \code{INT 21h}.
          Sertakan satu situasi ketika \code{T} lebih tepat dan satu
          situasi ketika \code{P} lebih tepat.

    \item Diberikan keluaran \code{D SS:00F4} berikut.

\begin{session}[caption={Keluaran \code{D} pada \eng{stack} untuk
soal uraian 2.}]
13A7:00F4  00 00 00 00 00 00 C3 1A-08 00 00 00 00 00 00 00  ................
\end{session}

          \begin{enumerate}
              \item Tentukan isi \reg{SS}:\hex{00FA} dan
                    \reg{SS}:\hex{00FC}.
              \item Bila \reg{SP} bernilai \hex{00FA} dan instruksi
                    berikutnya adalah \asm{RET}, alamat berapa yang
                    akan dimuat ke \reg{IP}? Jelaskan perhitungannya.
              \item Sebutkan satu penyebab yang paling mungkin bila
                    nilai tersebut bukan alamat kembali yang benar.
          \end{enumerate}

    \item Uraikan enam tahap metodologi \eng{debugging} sistematis.
          Untuk setiap tahap, sebutkan bukti yang harus dikumpulkan
          dan jelaskan mengapa tahap tersebut tidak boleh dilewati.

    \item Sebuah program menjumlahkan delapan byte larik, tetapi
          hasilnya selalu lebih besar dari yang seharusnya. Tuliskan
          langkah \eng{debugging} yang Anda lakukan mulai dari
          reproduksi sampai verifikasi. Sebutkan perintah DEBUG yang
          dipakai pada setiap langkah dan bukti yang diharapkan dari
          keluaran perintah tersebut.
\end{enumerate}

\subsection{Kisi-kisi Penilaian}

\begin{table}[H]
    \centering
    \caption{Kisi-kisi evaluasi Bab~\ref{chap:debugging}.}
    \label{tab:kisi-bab15}
    \small
    \begin{tabular}{@{}p{1.6cm}p{4.8cm}p{3.0cm}c@{}}
        \toprule
        \textbf{Sub-CPMK} & \textbf{Indikator yang diukur} &
        \textbf{Bentuk soal} & \textbf{Bobot} \\
        \midrule
        4.1 & Mengenali perintah penelusuran dan pemantauan \eng{flag} &
        Pilihan ganda 1, 3 & 15\% \\
        \addlinespace
        4.1 & Membedakan \code{T} dan \code{P} pada \asm{CALL} dan
        \code{INT} & Pilihan ganda 2; Uraian 1 & 25\% \\
        \addlinespace
        4.2 & Menganalisis \eng{dump} memori dan bingkai \eng{stack} &
        Uraian 2 & 25\% \\
        \addlinespace
        4.2 & Menjelaskan metodologi \eng{debugging} dan menerapkannya &
        Pilihan ganda 4, 5; Uraian 3, 4 & 35\% \\
        \bottomrule
    \end{tabular}
\end{table}

\begin{checklistbox}
\textbf{Checklist Pencapaian Kompetensi Bab~\ref{chap:debugging}}

Seorang mahasiswa dinyatakan kompeten pada Sub-CPMK 4.1 dan 4.2
apabila:
\begin{itemize}
    \item[\checked] dapat menyebutkan taksonomi kesalahan program
          beserta gejala khas masing-masing;
    \item[\checked] dapat memakai perintah \code{R}, \code{U},
          \code{T}, \code{P}, \code{G}, dan \code{D} tanpa panduan;
    \item[\checked] dapat menjelaskan kapan \code{T} dan kapan
          \code{P} dipakai pada \asm{CALL} serta \code{INT};
    \item[\checked] dapat mengisi tabel \eng{trace} register dan
          \eng{flag} dengan benar sekurang-kurangnya sepuluh langkah;
    \item[\checked] dapat membaca \eng{dump} memori dan menentukan
          isi bingkai \eng{stack};
    \item[\checked] dapat menemukan dan memperbaiki perulangan tak
          berujung, kesalahan \eng{off-by-one}, dan \eng{stack} tidak
          seimbang dalam batas waktu yang diberikan;
    \item[\checked] dapat menerapkan enam tahap metodologi
          \eng{debugging} dan menuliskan laporan \eng{trace} yang
          dapat diulang oleh orang lain.
\end{itemize}
\end{checklistbox}
